\subsection{统一优化框架}

\subsubsection{模型结构}

四个问题共用同一组决策变量：计划购电量 $b$、最终购电量 $q$、紧急购电量 $q^{em}$、充电量 $c$、放电量 $q^{dis}$ 与弃光量 $s$，并以时段末储电量 $E$ 为状态变量。以问题二为例，全年模型为

\begin{equation}
\begin{aligned}
\min_{b,\,q^{em},\,c,\,q^{dis},\,s,\,E}\quad
& C_{\text{total}}=\sum_{\tau=1}^{T}\Bigl(p_{\tau}b_{\tau}+\alpha_{em}p_{\tau}q^{em}_{\tau}\Bigr),
\quad \alpha_{em}=5,\quad T=52560
\end{aligned}
\end{equation}

\begin{equation}
\text{s.t.}\quad
b_{\tau}+q^{em}_{\tau}+PV_{\tau}\Delta t+q^{dis}_{\tau}=L_{\tau}\Delta t+c_{\tau}+s_{\tau},
\qquad 0\le s_{\tau}\le PV_{\tau}\Delta t
\end{equation}

\begin{equation}
0\le c_{\tau}\le 5000\Delta t=833.3333,\qquad
0\le q^{dis}_{\tau}\le 5000\Delta t\,\eta=750.0000
\end{equation}

\begin{equation}
E_{\tau}=E_{\tau-1}+\eta c_{\tau}-\frac{q^{dis}_{\tau}}{\eta},\qquad
\eta=0.9,\qquad 1200\le E_{\tau}\le 10800,\qquad E_{0}=6000
\end{equation}

\begin{equation}
b_{\tau}\ge0,\quad q^{em}_{\tau}\ge0,\quad c_{\tau}\ge0,\quad q^{dis}_{\tau}\ge0,\quad s_{\tau}\ge0
\end{equation}

式 (2) 的等式平衡使“供电不低于负载”自动成立，同时充电不与负载竞争；这一形式化与微网并网储能系统的逐时段成本最小建模传统一致\cite{ref-csee-2023-mg-ems-review,ref-moradi-2018-standalone,ref-alguhi-2025-bess-operation,ref-valibeygi-2020-battery-dispatch-soc}。模型规模为变量 $6T=315\,360$ 个、等式约束 $2T=105\,120$ 条、非零元 $9T-1=473\,039$ 个、整数变量 0 个；矩阵高度稀疏（非零元占比约 $1.4\times10^{-5}$），实现时采用稀疏存储，稠密存储需约 265 GB 内存并不可行。

\subsubsection{三处需要论证的建模要点}

\textbf{其一，是否需要二元变量强制“不同时充放电”。} 直观上储能不应在同一时段既充电又放电。本文不引入二元变量，而用如下引理保证线性松弛无损。

\begin{quote}
\textbf{引理 1（互补性）} 若 $\eta^2<1$ 且电价 $p_{\tau}>0$，则任一同时充放电的解都不是严格最优解；模型存在满足 $c_{\tau}q^{dis}_{\tau}=0$ 的最优解。

\textbf{证明} 设某时段 $c_{\tau}>0$、$q^{dis}_{\tau}>0$，令 $\varepsilon=\min(c_{\tau},q^{dis}_{\tau})>0$，把二者同时减去 $\varepsilon$：储电量净变化为 $-\eta\varepsilon+\varepsilon/\eta>0$（因 $\eta^2<1$），状态反而上升。因此可再把充电量减去 $\eta^2\varepsilon$，即取 $c'_{\tau}=c_{\tau}-\varepsilon$、$q^{dis\prime}_{\tau}=q^{dis}_{\tau}-\eta^2\varepsilon$，此时状态净变化恰为 $-\eta\varepsilon+\eta^2\varepsilon/\eta=0$，状态轨迹完全不变。该修改使购电量减少 $\varepsilon(1-\eta^2)>0$，费用严格下降，与最优性矛盾。故存在最优解满足互补性。$\square$
\end{quote}

引理 1 说明“不需要二元变量”是精确性质而非近似；其代价是最优面可能不唯一，故本文统一采用字典序规则选取提交解。同一现象在电池调度文献中亦被讨论：当往返效率小于 1 时线性松弛的互补性可以自发满足，而强制互补会带来数值与规模上的额外代价\cite{ref-wang-2024-complementarity,ref-nazir-2023-realizable-dispatch}。

\textbf{其二，光伏盈余如何闭合。} 附件 1 的光伏峰值 7612.32 kW 高于负载峰值 5958.97 kW，附件 2 中也有相当比例的时段光伏超过负载。若不加弃光变量而要求等式平衡，模型在这些时段将不可行；若改为不等式（净供给不低于负载），余电可以凭空消失，会低估可行域并扭曲费用。本文采用等式平衡加显式弃光变量（$s\in[0,PV\Delta t]$、目标系数为零）的做法：盈余只能被储能吸收或以 $s$ 弃掉，既不凭空消失也不上网售电。光伏限出力在含储能的经济调度中通常以同样的方式显式建模\cite{ref-grimaldi-2025-arbitrage-curtail,ref-qing-2025-pv-curtailment}。

\textbf{其三，终端条件。} 题面未规定决策期末的储电量，本文保留终端自由，只要求其落在运行区间内。由于期末库存具有经济价值，终端自由会使模型在末期“放空”储能，这一效应在结果中被显式量化（问题二为 $+2217.08$ 元），以免与价格套利混淆。

\subsubsection{结构恒等式}

对任意可行解，把储能状态转移式在决策期上求和可得两条与目标函数无关的恒等式，它们是比目标值更严格的正确性检验：

\begin{equation}
\sum_{\tau}q^{dis}_{\tau}=\eta^2\sum_{\tau}c_{\tau}-\eta\bigl(E_{T}-E_{0}\bigr)
\end{equation}

\begin{equation}
\sum_{\tau}b_{\tau}=N+\sum_{\tau}s_{\tau}+\bigl(1-\eta^2\bigr)\sum_{\tau}c_{\tau}
+\eta\bigl(E_{T}-E_{0}\bigr)-\sum_{\tau}q^{em}_{\tau},
\qquad N=\sum_{\tau}\bigl(L_{\tau}-PV_{\tau}\bigr)\Delta t
\end{equation}

式 (8) 的含义是：全年购电量等于净负荷加上弃光、加上储能循环损耗，再扣除期末放空的库存。以问题二为例，$\eta(E_T-E_0)=0.9\times(1200-6000)=-4320$ kWh，即期末放空等价于少购 4320 kWh；这条账在四个问题中都被独立复核，残差为零。

\subsubsection{四个维度上的参数化}

四个问题不需要分别推导模型，只需在表~\ref{tab:param} 的四个维度上取不同取值。

\begin{table}[htbp]
\centering
\footnotesize
\caption{四个问题在统一框架下的参数化}
\label{tab:param}
\begin{tabular}{p{1.4cm}p{2.4cm}p{3.4cm}p{3.6cm}p{2.6cm}}
\toprule
 & 时间维度 & 信息维度 & 结算维度 & 价格情景 \\
\midrule
问题一 & 单日 144 段 & 完全信息 & 按计划量计费 & 单日曲线（附件 1） \\
问题二 & 全年 52\,560 段 & 完全信息 & 计划量计费 + 5 倍紧急购电 & 逐日重复（附件 1） \\
问题三 & 全年，四个决策时刻 & 预报驱动（附件 3） & 上述 + 0.5/1.5 倍双向分段 & 逐日重复（附件 1） \\
链 \emph{4-2} & 全年，每日 0:00 & 负载/光伏完全信息，电价仅历史 & 计划量计费 + 5 倍紧急购电 & 实时波动，当日不可知 \\
链 \emph{4-3} & 全年，四个决策时刻 & 预报驱动，电价仅历史与已实现时段 & 上述 + 0.5/1.5 倍双向分段 & 实时波动，当日不可知 \\
\bottomrule
\end{tabular}
\end{table}

\subsubsection{求解与核验}

全部模型均为线性规划，采用 HiGHS 求解器在 CPU 上直接求解，不涉及 GPU：问题二的单次全年线性规划约 5.8 s，问题三与链 \emph{4-3} 因分解为 1460 个单日层分别约需 42 s 与 39 s，链 \emph{4-2} 的逐日递推约 12 s。所有计算在隔离任务、独立输出目录下由受监督进程执行，每次运行落盘解序列、交付表格、追踪链指纹与运行清单，可独立复现。

计算完成后对每个问题执行五项只读核验：产物与追踪链完整性（输入文件指纹未变、代码哈希与登记一致）；独立回代（用原始附件与解序列重建平衡方程、状态转移与约束，不读模型自带的自检结果、不重解）；量纲与逐层目标复算（含交付表格逐格比对）；解析界与常识检验；跨问题一致性（净负荷、弃光量等价格无关量逐位一致，而费用必然不同）。四问的等式残差均不超过 $4.6\times10^{-12}$，容量与功率约束零越界，交付工作簿逐格复算一致。
